\cvsection{Expériences professionnelles}

\begin{cventries}
  \cventry
    {Stagiaire ingénieur de recherche en IA}
    {INRIA}
    {Bordeaux, France}
    {Mars 2026 -- Septembre 2026}
    {
      \begin{cvitems}
        \item {Conception en autonomie, avec le suivi de mes encadrants, d’une application et d’un agent IA multi-étapes répondant au besoin d’aider les chercheurs à trouver des articles scientifiques.}
        \item {Développement de l’agent en Python sans framework d’orchestration pour comprendre ses mécanismes, avec API OpenAI, prompts adaptés aux outils et gestion de l’état et de la mémoire.}
        \item {Gestion du contexte avec sélection par l’utilisateur des messages de conversation à transmettre au LLM.}
        \item {Automatisation de la collecte de métadonnées d’articles scientifiques en ligne et transmission des résultats de tâches cron par webhooks.}
        \item {Construction d’un pipeline RAG pour exploiter les documents scientifiques : chunking, embeddings et recherche sémantique dans une base vectorielle ChromaDB.}
        \item {Développement d’API REST en Python/FastAPI documentées avec OpenAPI, de la persistance MongoDB et d’une interface React/Next.js pour interagir avec l’agent.}
        \item {Conteneurisation Docker, mise en place d’une CI/CD sur GitHub et gestion d’une application utilisant des LLM installés localement sur un Mac Studio.}
        \item {Utilisation de vLLM sur HPC avec GPU NVIDIA et de Slurm pour comparer Qwen, Gemma, GLM et plusieurs quantifications selon la qualité, la latence, le TTFT et la taille du contexte.}
        \item {Définition du protocole expérimental, analyse des résultats et rédaction en cours d’un article scientifique comparant les modèles.}
        \item {Évaluation d’OpenClaw dans un environnement Docker isolé et analyse des risques liés à l’accès autonome de l’agent aux outils et aux fichiers.}
      \end{cvitems}
    }

  \cventry
    {Stagiaire de recherche en apprentissage par renforcement}
    {Université Karunya}
    {Coimbatore, Inde}
    {Juin 2025 -- Sept. 2025}
    {
      \begin{cvitems}
        \item {Développement d’un modèle d’apprentissage par renforcement et de son environnement d’entraînement pour la navigation autonome d’un robot hospitalier.}
        \item {Déploiement et optimisation du logiciel et du modèle sur Raspberry Pi pour des tests en conditions réelles.}
      \end{cvitems}
    }
\end{cventries}
